\documentclass{exam}

\usepackage[margin=1in]{geometry}

\usepackage{comment}
\usepackage{graphicx}
\usepackage{listings}

\pagestyle{headandfoot}
\firstpageheader{CSCI 221}{}{Pointers and Arrays Assignment}
\runningheader{CSCI 221}{}{Pointers and Arrays Assignment}


\begin{document}
\title{Pointers and Arrays Assignment (Evaluation)}
\author{CSCI 221: Computer Science Fundamentals 2}
\date{Spring 2026}
\maketitle


This assignment is designed to test your understanding of pointers and arrays. 
Feel free to collaborate with others and use resources, but all code and writeups must be your own. 
To submit your assignment, please submit your code, makefile, and writeup. 
For full credit, your code should check for invalid input. 

Some questions in this assignment allow you to add additional arguments to your functions. 
For full credit, you should only do this when necessary. 

\textbf{Due Date:}
Friday, February 27th at 1:00 pm. 

\begin{questions}

\question[30]
\textbf{Binary Search Trees.}
In this question, you will implement a binary search tree data structure. 
A binary search tree contains nodes that store three items: a data value, a pointer to the left child node, and a pointer to the right child node. 
Any pointer to a child that does not exist is set to \texttt{NULL}. 
For this question, we will design a tree that stores strings, but does not allow duplicate values. 

Binary search trees are designed such that all values in the left subtree of a node must be less than the value in the node, and all values in the right subtree of a node must be greater than the value in the node. 
A subtree of a node is recursively defined as the tree consisting of the node and its subtrees, rooted at the associated child of that node. 
That is, all nodes that can be reached by following the pointer to the left child of node A are in node A's left subtree. 

In order to create a tree of strings, we must define what less than and greater than mean on strings. 
For this question our comparison will compare the 0th characters of the strings. 
If they are different, we say that the relationship between strings is the same as that between those characters. 
If they are equal, we move to the next character and repeat the process. 
Because we use the \texttt{NUL} character and do not allow duplicates, we should never have to compare an existing character to a non-existing character. 

For this question, you are not allowed to use functions in the \texttt{<string>} library. 

\begin{parts}
\part[3]
Design a struct to store a node in the binary search tree. 

\part[2]
Design a struct to represent a binary search tree. 
This struct should contain a pointer to the root node in the tree and the number of elements in the tree. 
Remember to update the count in other functions as appropriate. 

\part[5]
Write a function to compare two strings according to the process above. 
Your function should take as input two strings. 
It should return an enum representing the relationship between those strings. 
You are NOT allowed to add additional arguments to your function. 

\part[5]
Write a function that takes as input a binary search tree and a string and adds a copy of the string to the binary search tree. 
Make sure to add a deep copy rather than a shallow copy. 
You should follow the BST property to ensure that you add the new node to the correct place in the tree. 
This should not change any other part of the tree. 
The function should have no return value. 

\part[10]
Write a function that takes as input a binary search tree and a value and removes the value from the tree. 
The function should modify the original tree, not create a copy. 
The tree must be updated to maintain its structure. 
If possible, the left child should be moved to the position occupied by the removed value. 
This may cause additional issues that need to be fixed. 
The function should return the number of values removed (0 or 1). 
Beware of memory leaks. 

\part[5]
Write a function that deallocates a binary search tree. 
The function should take as input a pointer to a binary search tree and should deallocate the tree and all nodes contained within it. 
\end{parts}

\question[5]
\textbf{I/O.}
Write a main function that prints out copies of user input. 
You should ask the user for a string and how many copies of it they want made. 
Then, you should print out that number of copies of that string on separate lines. 
However, if your main function is passed the \texttt{-r} flag when it is executed, it should print out each copy of the user input in reverse character order. 
For example, giving the function \texttt{"Hello"} and \texttt{2} with the \texttt{-r} flag would result in:
\begin{verbatim}
olleh
olleh
\end{verbatim}
Your function should only respond to the \texttt{-r} flag, but it should accept it with any other arguments. 
Those other arguments should be ignored. 

\question[10]
\textbf{Writeup.}
In addition to your code, you should submit a brief (~1-2 pages) writeup that describes the state of your code. 
You should discuss what is working, what is not, and any known errors. 
In this discussion, you should reference the tests you did and how they support your claims. 
You will need to provide a brief description of your tests for you to reference. 
Your writeup should be in a separate .txt or .pdf file. 

\end{questions}

\end{document}